\chapter{Glossary}
\label{ch:glossary}

The glossary uses the meanings developed in this book. Many terms have wider
or competing meanings in the literature; the definitions below are working
definitions for \RCR.

\begin{description}[style=nextline,leftmargin=0pt,labelindent=0pt,labelwidth=0pt]
\item[Abduction.] Inference that selects a hypothesis because it explains the
evidence better than relevant alternatives. It is ampliative, not deductively
certain.
\index{abduction}

\item[Accuracy.] Closeness of a measurement or representation to the target
value or structure. Accuracy differs from precision.
\index{accuracy}

\item[Ad hoc adjustment.] A modification introduced mainly to protect a claim
from known counterevidence without generating independent new constraints.
\index{ad hoc adjustment}

\item[Bayesian updating.] Revision of probabilities using Bayes' theorem:
posterior probability is proportional to likelihood times prior probability.
\index{Bayesian updating}

\item[Calibration.] Establishing how an instrument's response relates to a
standard or independently measured quantity.
\index{calibration}

\item[Causal effect.] A difference between outcomes under specified
interventions or potential treatments, for a defined unit or population.
\index{causal effect}

\item[Causal mechanism.] Organized parts, activities, and relations whose
interaction produces a phenomenon.
\index{causal mechanism}

\item[Claim.] A proposition with a target, content, scope, alternatives, and
warrant. \RCR\ treats a claim as incomplete if these fields remain hidden.
\index{claim}

\item[Construct validity.] The degree to which an operation measures the
construct it is intended to represent.
\index{construct validity}

\item[Constraint.] A stable restriction or tendency on what can happen under
specified conditions. \RCR\ distinguishes material, formal, causal,
organizational, and institutional constraints.
\index{constraint}

\item[Constructive empiricism.] Van Fraassen's view that science aims at
empirical adequacy and need not require belief in the truth of unobservable
theoretical claims.
\index{constructive empiricism}

\item[Counterfactual.] A conditional claim about what would happen under a
condition that may not actually occur.
\index{counterfactual}

\item[Correction capacity.] The ability of an inquiry system to expose, assess,
and repair error.
\index{correction capacity}

\item[Data.] Local records produced by measurement, observation, or collection.
Data become evidence for phenomena through analysis and validation.
\index{data}

\item[Demarcation.] The problem of distinguishing science, non-science, and
pseudoscience. \RCR\ treats it as a profile of epistemic performance rather
than a single binary property.
\index{demarcation}

\item[Emergence.] The occurrence of stable, explanatory, or causal patterns at a
level of organization that depends on, but is not simply interchangeable with,
lower-level description.
\index{emergence}

\item[Empirical adequacy.] Correct description of observable phenomena within a
specified scope, whether or not a theory's unobservable posits are accepted.
\index{empirical adequacy}

\item[Entity realism.] The position that successful intervention with a
theoretical entity can support belief in the entity's reality.
\index{entity realism}

\item[Epistemic uncertainty.] Uncertainty caused by limited knowledge of
parameters, models, initial conditions, or measurement processes.
\index{epistemic uncertainty}

\item[Error correlation.] The extent to which apparently independent methods
share a source of failure.
\index{error correlation}

\item[Error-statistical reasoning.] An approach that emphasizes what a test
would probably have detected if an error or discrepancy were present.
\index{error-statistics}

\item[Explanation.] An answer to a why, how, or under-what-conditions question.
Explanation may be deductive, unificatory, causal, mechanistic, or model-based.
\index{explanation}

\item[Falsificationism.] Popper's emphasis on risky claims and attempts at
refutation as central to scientific rationality.
\index{falsificationism}

\item[Fallibilism.] The view that any finite belief or claim may require
revision, even when well supported.
\index{fallibilism}

\item[Inference to the best explanation.] Comparative reasoning that favours a
hypothesis because it explains evidence better than relevant alternatives.
\index{inference to the best explanation}

\item[Incommensurability.] Difficulty of direct translation or comparison
between conceptual frameworks, especially across major theory change.
\index{incommensurability}

\item[Invariance.] Stability of a relevant relation under specified changes in
observer, instrument, representation, intervention, or context.
\index{invariance}

\item[Intervention boundary.] The variables changed, processes allowed to
respond, and background conditions held fixed in a causal claim.
\index{intervention boundary}

\item[Idealization.] Deliberate simplification that removes or alters features
of a target in order to investigate a selected relation.
\index{idealization}

\item[Law-like constraint.] A stable modal relation that supports counterfactual
or interventional reasoning within a defined regime.
\index{laws of nature}

\item[Likelihood.] The probability of evidence conditional on a hypothesis,
written \(\Prob(E\mid H)\). It is not the probability that the hypothesis is
true.
\index{likelihood}

\item[Measurement path.] The chain from target through interaction, instrument,
calibration, transformation, error model, and interpretation.
\index{measurement path}

\item[Model.] A selective representation used to reason about a target. Models
may be mathematical, computational, diagrammatic, mechanistic, statistical,
or material.
\index{model}

\item[Normal science.] Kuhn's term for puzzle-solving research conducted within
a shared paradigm.
\index{normal science}

\item[Objectivity.] In \RCR, the graded achievement of world-sensitive,
publicly corrigible, scope-honest inquiry.
\index{objectivity}

\item[Ontological commitment.] The degree and kind of belief that evidence
warrants about entities, structures, processes, or social facts.
\index{ontological commitment}

\item[Paradigm.] A historically situated framework of exemplars, concepts,
standards, instruments, and problems that organizes research.
\index{paradigm}

\item[Pragmatic adequacy.] The degree to which a representation supports
reliable inquiry, prediction, control, diagnosis, or design.
\index{pragmatic adequacy}

\item[Precision.] The spread of repeated measurements. Precision does not imply
accuracy.
\index{precision}

\item[Prior probability.] A probability assigned to a hypothesis before a
specified item of evidence is incorporated.
\index{prior probability}

\item[Public corrigibility.] The institutional capacity for others to inspect,
criticize, reproduce, and revise a claim.
\index{public corrigibility}

\item[Reflexive objectivity.] Objectivity that includes explicit awareness and
critique of the social, conceptual, and material conditions of inquiry.
\index{reflexive objectivity}

\item[Reflexive Constraint Realism.] The framework developed in this book:
mind-independent generative reality, layered representation, graded commitment,
and public correction through constraint-sensitive inquiry.
\index{Reflexive Constraint Realism}

\item[Replication.] Repeating or extending an inquiry under a specified relation
of similarity in order to test stability, error, or transport.
\index{replication}

\item[Research programme.] A historically extended set of theoretical
commitments, auxiliary hypotheses, heuristics, and problem-solving practices.
\index{research programme}

\item[Scope condition.] A statement of the objects, population, time, scale,
interventions, or background circumstances to which a claim applies.
\index{scope condition}

\item[Scientific realism.] The view that successful scientific theories can
warrant belief in both observable and some unobservable aspects of reality.
\index{scientific realism}

\item[Situated knowledge.] Knowledge produced from a partial social and material
location whose limits should be made explicit.
\index{situated knowledge}

\item[Standpoint theory.] The view that social positions can shape what is
visible and can sometimes provide epistemic advantages concerning relations of
power, subject to critical testing.
\index{standpoint theory}

\item[Structural realism.] A family of views emphasizing that scientific
knowledge or reality is fundamentally relational or structural.
\index{structural realism}

\item[Theory-ladenness.] Dependence of observation, measurement, or
interpretation on prior concepts, theories, training, and instruments.
\index{theory-ladenness}

\item[Transportability.] The validity of carrying a result from one setting to
another after checking shared mechanisms and changed conditions.
\index{transportability}

\item[Truth.] In \RCR, a layered relation involving correspondence, structural
accuracy, interventional adequacy, and pragmatic adequacy. These layers may
diverge.
\index{truth}

\item[Value transparency.] Explicit reporting of value choices that shape
questions, concepts, error thresholds, institutional design, or action.
\index{value transparency}

\item[World resistance.] The fact that some claims fail because independent
processes do not behave as the claim requires.
\index{world resistance}
\end{description}

\section*{How to use the glossary}

The glossary is not a substitute for the arguments. In particular, the terms
\emph{constraint}, \emph{objectivity}, and \emph{truth} acquire their specific
meaning through the relations among the axioms, inferential rules, and case
analyses. A reader should return to Chapters~\ref{ch:ontology} and
\ref{ch:epistemology} when a definition appears to conflict with a familiar use.
